\clearpage
\subsection{Die natürliche Exponentialfunktion}
\label{subsec:real_functions_natural_exponential}
Eine besondere Basis ergibt sich, wenn ein Gesamtzuwachs von $100\,\%$
in immer mehr gleich große, aufeinander aufbauende Schritte zerlegt wird.
\begin{center}
  \renewcommand{\arraystretch}{1.4}
  \begin{tabular}{@{}c|rrrr@{}}
    Anzahl $n$&1&10&100&1000\\
    $(1+1/n)^n$&$2$&$2{,}5937$&$2{,}7048$&$2{,}7169$
  \end{tabular}
\end{center}
\begin{DefinitionBox}{Eulersche Zahl und natürliche Exponentialfunktion}{real_functions_euler_number}
  Die \emph{Eulersche Zahl} \EnglishTerm{Euler's number} ist
  \[
    e=\lim_{n\to\infty}\left(1+\frac1n\right)^n
      =2{,}7182818\dots
  \]
  Die Funktion $\exp(x)=e^x$ heißt
  \emph{natürliche Exponentialfunktion}
  \EnglishTerm{natural exponential function}.
\end{DefinitionBox}
Die Tabelle motiviert den Grenzwert; seine Existenz wird hier nicht
bewiesen. Die besondere Bedeutung der Basis $e$ wird bei der
Differentialrechnung deutlich.

\begin{ExampleBox}{Abklingen einer Schwingungshülle}{real_functions_exponential_decay}
  Die Amplitude einer gedämpften Schwingung werde durch
  $A(t)=8\,\mathrm{mm}\cdot e^{-t/\tau}$ für $t\ge0$ mit
  $\tau=2\,\mathrm{s}$ beschrieben.
  Bestimmen Sie den Anfangswert, den Wert nach einer Zeitkonstante
  und das langfristige Verhalten.
  \begin{ClozeArea}[34]{7}
    Der Exponent ist dimensionslos. Aus den Eigenschaften von $e^x$ folgt
    \[
      \begin{aligned}
        A(0)&=8\,\mathrm{mm},\\
        A(\tau)&=\frac8e\,\mathrm{mm}\approx2{,}94\,\mathrm{mm},\\
        \lim_{t\to+\infty}A(t)&=0.
      \end{aligned}
    \]
    Die Funktion ist auf $[0,\infty)$ streng monoton fallend und besitzt
    die \textcolor{green!55!black}{Wertemenge} $(0,8\,\mathrm{mm}]$.
  \end{ClozeArea}
\end{ExampleBox}

\begin{KeyPointBox}{Exponentielle Modelle lesen}{real_functions_exponential_model}
  Bei $A(t)=A_0e^{kt}$ mit $A_0>0$ beschreibt $A_0$ den Anfangswert.
  Das Vorzeichen von $k$ entscheidet:
  $k>0$ bedeutet Wachstum, $k<0$ Abnahme, $k=0$ einen konstanten Wert.

  Ist $t$ eine Zeit, hat $k$ die Einheit einer inversen Zeit.
  Der Exponent $kt$ muss dimensionslos sein.
\end{KeyPointBox}
Wie lange es dauert, bis ein bestimmter Wert erreicht wird, können wir
noch nicht allgemein berechnen. Dafür benötigen wir später die
Umkehrung der Exponentialfunktion: den Logarithmus.
